% Partie : méthodes
% Chapitre : démonstrations
% Section : lire des démonstrations

\documentclass[../../../include/open-logic-section]{subfiles}

\begin{document}
\olfileid[fr]{mth}{prf}{rea}
\olsection{Lire des démonstrations}

Dans les manuels et les articles, les
démonstrations donnent rarement autant
de détails que les exemples étudiés
jusqu'ici. Les auteurs ne signalent
pas toujours qu'ils distinguent des
cas, qu'ils procèdent indirectement
ou qu'ils appliquent une définition.
Lorsque vous lisez une démonstration,
vous devez donc souvent reconstituer
vous-même ces étapes pour la comprendre.
C'est aussi un bon exercice pour apprendre à reconnaître les étapes
d'un raisonnement.
Examinons un exemple.

\begin{prop}[Absorption]
Pour tous ensembles $A$ et $B$,
\[
A \cap (A \cup B) = A.
\]
\end{prop}

\begin{proof}
Si $z \in A \cap (A \cup B)$,
alors $z \in A$ ; ainsi,
$A \cap (A \cup B) \subseteq A$.
Supposons maintenant $z \in A$.
Alors $z \in A \cup B$, donc
$z \in A \cap (A \cup B)$.
\end{proof}

Cette démonstration de la loi
d'absorption est très condensée.
Elle ne nomme pas les définitions
utilisées et ne dit pas «~nous
devons montrer que…~» avant de
démontrer chaque étape. Déplions-la.
La proposition concerne tous
ensembles $A$ et $B$ ; dans la
démonstration, ces lettres désignent
des ensembles arbitraires.
Les énoncés généraux qu'elle établit
sont exactement ceux qu'exige
l'égalité des ensembles : tout
!!{element} du membre gauche est
!!a{element} du membre droit,
et réciproquement.

\begin{quote}
«~Si $z \in A \cap (A \cup B)$,
alors $z \in A$ ; ainsi,
$A \cap (A \cup B) \subseteq A$.~»
\end{quote}

Voilà la première moitié de la
démonstration : si un~$z$ arbitraire
est !!a{element} du membre gauche,
il est aussi !!a{element} du membre
droit, d'où
$A \cap (A \cup B) \subseteq A$.
Supposons $z \in A \cap (A \cup B)$.
Un objet est !!{element} de
l'intersection de deux ensembles
si et seulement s'il est
!!{element} de chacun
d'eux. Nous pouvons donc conclure
que $z \in A$ et $z \in A \cup B$.
En particulier, $z \in A$,
comme il fallait le montrer.
La première moitié étant achevée,
le reste doit démontrer
$A \subseteq A \cap (A \cup B)$.

\begin{quote}
«~Supposons maintenant $z \in A$.
Alors $z \in A \cup B$, donc
$z \in A \cap (A \cup B)$.~»
\end{quote}

Nous commençons par supposer
$z \in A$, car il faut montrer
que, pour tout~$z$, si
$z \in A$, alors
$z \in A \cap (A \cup B)$.
Pour établir cette dernière
appartenance, la définition
de~$\cap$ demande
(i) $z \in A$ et
(ii) $z \in A \cup B$.
La condition (i) est notre
hypothèse : il n'y a rien
à ajouter, ce qui explique
que la démonstration ne la
mentionne plus. Pour (ii),
rappelons qu'un objet est
!!a{element} d'une union
si et seulement s'il est
!!{element} d'au moins un
des ensembles réunis.
Comme $z \in A$ et que
$A \cup B$ est l'union
de $A$ et de $B$, nous
avons bien $z \in A \cup B$.
Les conditions (i) et (ii)
sont établies ; par définition
de~$\cap$,
$z \in A \cap (A \cup B)$.
La démonstration initiale
ne cite aucune de ces
définitions : elle suppose
que le lecteur les a déjà
assimilées. Sinon, il faut
y revenir, puis reconnaître
pourquoi $z \in A$ entraîne
$z \in A \cup B$ et pourquoi
$z \in A$ avec
$z \in A \cup B$ entraîne
$z \in A \cap (A \cup B)$.

Voici une version de la même
démonstration où toutes les
étapes sont explicitées :
\begin{proof}{}
[Par définition de l'égalité des
ensembles, pour démontrer
$A \cap (A \cup B) = A$,
il faut établir
(a) $A \cap (A \cup B) \subseteq A$
et
(b) $A \subseteq A \cap (A \cup B)$.
Pour (a), la définition de
$\subseteq$ nous demande
de montrer que
$z \in A \cap (A \cup B)$
implique $z \in A$.]
Si $z \in A \cap (A \cup B)$,
alors $z \in A$
[car, par définition de
$\cap$, $z$ appartient à
$A \cap (A \cup B)$
si et seulement s'il
appartient à $A$ et
à $A \cup B$] ;
donc
$A \cap (A \cup B) \subseteq A$.
[(b) : par définition de
$\subseteq$, il faut
montrer que $z \in A$
implique
$z \in A \cap (A \cup B)$.]
Supposons maintenant
[(1)] $z \in A$.
Alors
[(2)] $z \in A \cup B$
[puisque, d'après (1),
$z \in A$ ou $z \in B$,
ce qui signifie
$z \in A \cup B$
par définition de~$\cup$] ;
par conséquent,
$z \in A \cap (A \cup B)$
[car la définition
de~$\cap$ exige
$z \in A$, c'est-à-dire
(1), et
$z \in A \cup B$,
c'est-à-dire (2)].
\end{proof}

\begin{prob}
Développez la démonstration
suivante de
$A \cup (A \cap B) = A$.
Mentionnez tous les schémas
d'inférence utilisés,
expliquez pourquoi chaque
étape découle des hypothèses
ou des conclusions déjà
établies, et précisez
les définitions invoquées.
\begin{proof}
  Si $z \in A \cup (A \cap B)$,
  alors $z \in A$ ou
  $z \in A \cap B$.
  Si $z \in A \cap B$,
  alors $z \in A$.
  Tout $z \in A$ appartient
  aussi à $A \cup (A \cap B)$.
\end{proof}
\end{prob}

\end{document}
